% Options for packages loaded elsewhere
\PassOptionsToPackage{unicode}{hyperref}
\PassOptionsToPackage{hyphens}{url}
\documentclass[
]{article}
\usepackage{xcolor}
\usepackage[margin=1in]{geometry}
\usepackage{amsmath,amssymb}
\setcounter{secnumdepth}{-\maxdimen} % remove section numbering
\usepackage{iftex}
\ifPDFTeX
  \usepackage[T1]{fontenc}
  \usepackage[utf8]{inputenc}
  \usepackage{textcomp} % provide euro and other symbols
\else % if luatex or xetex
  \usepackage{unicode-math} % this also loads fontspec
  \defaultfontfeatures{Scale=MatchLowercase}
  \defaultfontfeatures[\rmfamily]{Ligatures=TeX,Scale=1}
\fi
\usepackage{lmodern}
\ifPDFTeX\else
  % xetex/luatex font selection
\fi
% Use upquote if available, for straight quotes in verbatim environments
\IfFileExists{upquote.sty}{\usepackage{upquote}}{}
\IfFileExists{microtype.sty}{% use microtype if available
  \usepackage[]{microtype}
  \UseMicrotypeSet[protrusion]{basicmath} % disable protrusion for tt fonts
}{}
\makeatletter
\@ifundefined{KOMAClassName}{% if non-KOMA class
  \IfFileExists{parskip.sty}{%
    \usepackage{parskip}
  }{% else
    \setlength{\parindent}{0pt}
    \setlength{\parskip}{6pt plus 2pt minus 1pt}}
}{% if KOMA class
  \KOMAoptions{parskip=half}}
\makeatother
\usepackage{graphicx}
\makeatletter
\newsavebox\pandoc@box
\newcommand*\pandocbounded[1]{% scales image to fit in text height/width
  \sbox\pandoc@box{#1}%
  \Gscale@div\@tempa{\textheight}{\dimexpr\ht\pandoc@box+\dp\pandoc@box\relax}%
  \Gscale@div\@tempb{\linewidth}{\wd\pandoc@box}%
  \ifdim\@tempb\p@<\@tempa\p@\let\@tempa\@tempb\fi% select the smaller of both
  \ifdim\@tempa\p@<\p@\scalebox{\@tempa}{\usebox\pandoc@box}%
  \else\usebox{\pandoc@box}%
  \fi%
}
% Set default figure placement to htbp
\def\fps@figure{htbp}
\makeatother
\setlength{\emergencystretch}{3em} % prevent overfull lines
\providecommand{\tightlist}{%
  \setlength{\itemsep}{0pt}\setlength{\parskip}{0pt}}
\usepackage[spanish]{babel}
\usepackage{amsmath}
\usepackage{fontspec}
\usepackage{unicode-math}
\usepackage{graphicx}
\usepackage{adjustbox}
\usepackage{tikz}
\usepackage{pgfplots}
\usepackage{booktabs}
\usetikzlibrary{babel}
\pgfplotsset{compat=1.18}
\usepackage{bookmark}
\IfFileExists{xurl.sty}{\usepackage{xurl}}{} % add URL line breaks if available
\urlstyle{same}
\hypersetup{
  hidelinks,
  pdfcreator={LaTeX via pandoc}}

\author{}
\date{\vspace{-2.5em}}

\begin{document}

\section{Question}\label{question}

Una vez inicia la migración del bonito en la costa del Pacífico, su
pesca varía diariamente según la expresión:

\[P = -d^2 + 8d + -15\]

Donde \(P\) representa la pesca, en toneladas, y \(d\) los días
transcurridos después de iniciada la migración.

¿Cuál de las siguientes gráficas muestra correctamente la relación entre
la pesca y los días transcurridos después de iniciada la migración?

\subsection{Answerlist}\label{answerlist}

\begin{itemize}
\item
  \includegraphics[width=0.6\linewidth,height=\textheight,keepaspectratio]{grafica_opcion_a.png}
\item
  \includegraphics[width=0.6\linewidth,height=\textheight,keepaspectratio]{grafica_opcion_b.png}
\item
  \includegraphics[width=0.6\linewidth,height=\textheight,keepaspectratio]{grafica_opcion_c.png}
\item
  \includegraphics[width=0.6\linewidth,height=\textheight,keepaspectratio]{grafica_opcion_d.png}
\end{itemize}

\section{Solution}\label{solution}

Para resolver este problema, debemos identificar cuál gráfica representa
correctamente la función cuadrática \(P = -d^2 + 8d + -15\).

\textbf{Análisis de la función:}

\begin{enumerate}
\def\labelenumi{\arabic{enumi}.}
\tightlist
\item
  \textbf{Tipo de función}: Es una función cuadrática (parabólica)
\item
  \textbf{Coeficiente principal}: \(a = -1\) (negativo), por lo que la
  parábola \textbf{abre hacia abajo} (forma de U invertida)
\item
  \textbf{Vértice (máximo)}: Ocurre en
  \(d = \frac{-b}{2a} = \frac{8}{2} = 4\) días
\item
  \textbf{Valor máximo}: \(P_{max} = 1\) toneladas
\end{enumerate}

\textbf{Características que debe tener la gráfica correcta:}

\begin{itemize}
\tightlist
\item
  Parábola invertida (forma de U invertida)
\item
  Máximo alrededor del día 4
\item
  Los valores de pesca aumentan hasta el día del vértice y luego
  disminuyen
\end{itemize}

\textbf{Verificación con puntos calculados:}

Para \(d = 5\): \(P = -(5)^2 + 8(5) + -15 = -25 + 40 + -15 = 0\)
toneladas

\subsection{Answerlist}\label{answerlist-1}

\begin{itemize}
\tightlist
\item
  \textbf{Falso}. Esta gráfica muestra un patrón cuadrático ascendente
  (puntos que aumentan continuamente), lo cual no corresponde a una
  parábola invertida con coeficiente negativo \(-d^2\).
\item
  \textbf{Verdadero}. Esta gráfica muestra correctamente una parábola
  invertida (forma de U invertida) con máximo alrededor del día 4,
  conforme a la función \(P = -d^2 + 8d + -15\). Los valores aumentan
  hasta el vértice y luego disminuyen.
\item
  \textbf{Falso}. Esta gráfica muestra un patrón lineal descendente
  (línea recta decreciente), no una función cuadrática parabólica.
\item
  \textbf{Falso}. Esta gráfica muestra una parábola en forma de U que
  abre hacia arriba (valores disminuyen al centro y aumentan en
  extremos), contrario al comportamiento esperado con coeficiente
  negativo \(-d^2\).
\end{itemize}

\section{Meta-information}\label{meta-information}

exname: Migración Atún Representación Gráfica Función Cuadrática extype:
schoice exsolution: 0100 exshuffle: TRUE exsection:
Estadística/Interpretación de Gráficas/Funciones Cuadráticas

\end{document}
